\documentclass[11pt,a4paper]{article}
\usepackage[margin=1in]{geometry}
\usepackage[spanish]{babel}
\usepackage[utf8]{inputenc}
\usepackage[T1]{fontenc}
\usepackage{array}
\usepackage{booktabs}
\usepackage{enumitem}
\usepackage{graphicx}
\usepackage{hyperref}
\usepackage{siunitx}
\usepackage{tabularx}
\setlength{\parskip}{6pt}
\setlength{\parindent}{0pt}
\sisetup{output-decimal-marker={,}}

\title{An\'alisis V3 del l\'imite IMU-only y cierre del tramo activo}
\author{MOOVIA Sensor Lab}
\date{\today}

\begin{document}
\maketitle

\section*{Resumen ejecutivo}
Los JSON V3 confirman un cambio importante en el diagn\'ostico del proyecto: el cuello principal ya no es la transmisi\'on BLE. En las dos sesiones analizadas no aparecen p\'erdidas reales por \texttt{seq16} (\texttt{missingPackets = 0}), pero la trayectoria segu\'ia saliendo desvirtuada porque el algoritmo prolongaba el tramo activo cuando el gesto real ya hab\'ia terminado y la velocidad integrada quedaba contaminada.

En esta iteraci\'on se implementa una separaci\'on expl\'icita entre:
\begin{itemize}[leftmargin=14pt]
  \item \textbf{endpoint activo}: \'ultimo punto del gesto tras el recorte y estabilizaci\'on del tramo activo;
  \item \textbf{endpoint asentado}: posici\'on media del reposo final.
\end{itemize}

El resultado no convierte el sistema en un ``flight log'' absoluto. S\'i mejora la descripci\'on del gesto, deja visible el l\'imite real del algoritmo IMU-only y prepara la base para una Iteraci\'on 2 con sensores auxiliares.

\section{Contexto}
Las pruebas V3 corresponden a gestos estructurados de barra con el sensor aproximadamente mirando hacia arriba:
\begin{itemize}[leftmargin=14pt]
  \item \textbf{suelo $\rightarrow$ mesa}: desplazamiento vertical de referencia aproximado de \SI{1}{m};
  \item \textbf{mesa $\rightarrow$ suelo}: desplazamiento vertical de referencia aproximado de \SI{-1}{m};
  \item la captura empieza al pulsar \emph{grabar} y termina al pulsar \emph{parar}, por lo que incluye tiempo de colocaci\'on antes del gesto y reposo final despu\'es del gesto.
\end{itemize}

Ese detalle es clave: si todo el JSON se integra como si fuera un movimiento continuo, la trayectoria queda sesgada por la cola quieta y por el ruido integrado en reposo.

\section{Estado previo observado en V3}
Antes del cambio actual, los V3 ya mostraban una captura sana por radio, pero un cierre deficiente del tramo activo:
\begin{itemize}[leftmargin=14pt]
  \item \texttt{V3-moovia-session-1774637821525.json}: \texttt{missingPackets=0}, \texttt{avgRate\,$\approx$\,1060,8 Hz}, pero el tramo activo terminaba con \(v_z \approx \SI{-1,846}{m/s}\), \(z \approx \SI{-1,674}{m}\), \(|a_{lin}| \approx \SI{0,127}{m/s^2}\) y gyro \(\approx \SI{0,077}{dps}\). El gesto real ya hab\'ia terminado, pero la integraci\'on segu\'ia ``corriendo''.
  \item \texttt{V3-moovia-session-1774637899200.json}: tambi\'en con \texttt{missingPackets=0} y una forma vertical coherente (``sube un poco antes de bajar''), pero el tramo activo a\'un acababa con \(v_z \approx \SI{-0,634}{m/s}\) y \(z \approx \SI{-0,336}{m}\).
\end{itemize}

La conclusi\'on t\'ecnica era clara: la radio ya no explicaba la deformaci\'on principal. El error dominante estaba en el cierre tard\'io del movimiento, la velocidad residual no anulada, la sensibilidad del yaw/lateral y la posible fijaci\'on demasiado temprana del eje de barra (\texttt{barAxis}).

\section{Cambios implementados}
\subsection*{1. Endpoint activo frente a endpoint asentado}
El an\'alisis compartido en \texttt{sensor-core} ahora distingue entre:
\begin{itemize}[leftmargin=14pt]
  \item \textbf{Active end point}: \'ultimo punto del tramo activo mostrado;
  \item \textbf{Settled end point}: posici\'on media del reposo final.
\end{itemize}

Con esto las gr\'aficas y la visualizaci\'on 3D ense\~nan la reconstrucci\'on activa estabilizada, mientras que la posici\'on asentada final se reporta aparte.

\subsection*{2. Recorte de cola quieta y ``velocity kill'' de postprocesado}
Cuando al final de la sesi\'on aparece una ventana suficientemente larga con baja energ\'ia f\'isica (\(|a_{lin}|\) bajo y gyro bajo), el algoritmo:
\begin{enumerate}[leftmargin=16pt]
  \item recorta esa cola del tramo activo;
  \item fija el motivo de cierre como \texttt{quiet\_tail};
  \item fuerza la velocidad residual reportada del tramo activo a cero en la reconstrucci\'on estabilizada.
\end{enumerate}

Esto no modifica el JSON crudo exportado; solo afecta al an\'alisis offline y a la visualizaci\'on.

\subsection*{3. `barAxis' solo con giro real}
La estimaci\'on del eje de manga ya no se congela durante el reposo inicial. Solo se acumulan muestras de gyro suficientemente energ\'eticas y se reporta una confianza:
\begin{itemize}[leftmargin=14pt]
  \item \texttt{high}: eje dominante consistente;
  \item \texttt{low}: evidencia insuficiente o poco dominante;
  \item \texttt{unavailable}: no hay datos suficientes para estimarlo.
\end{itemize}

Si no hay evidencia suficiente, el sistema degrada a ``sin compensaci\'on de manga'' en vez de inventar un eje.

\subsection*{4. Diagn\'ostico temporal}
Se a\~nade una m\'etrica diagn\'ostica de tiempo:
\begin{itemize}[leftmargin=14pt]
  \item \textbf{\texttt{effectiveTickUs} de captura}: inferido a partir de \texttt{hwTs16}, alrededor de \SI{1,067}{\micro\second}/tick;
  \item \textbf{\texttt{effectiveTickUs} del an\'alisis}: mediana del intervalo reconstruido entre muestras, alrededor de \SI{999,5}{\micro\second}, es decir, muy cerca de \SI{1}{ms}.
\end{itemize}

Esto sugiere que el \(\approx \SI{1060}{Hz}\) observado no equivale necesariamente a una ODR real de \SI{1060}{Hz}; parte de esa cifra viene de la interpolaci\'on temporal y de intervalos sub-milisegundo dentro del pipeline.

\section{Estado actual de V3 tras el cambio}
La tabla~\ref{tab:v3-current} resume el estado actual de las dos sesiones V3 tras aplicar el nuevo postprocesado.

\begin{table}[h!]
\centering
\small
\resizebox{\textwidth}{!}{%
\begin{tabular}{lrrrrrrrrrr}
\toprule
Archivo & Rate [Hz] & missPkt & missSmpl & maxGap [ms] & Activo [s] & Pico lin [m/s\textsuperscript{2}] & Alt. m\'ax [m] & Alt. final [m] & Lat. m\'ax [m] & Lat. final [m] \\
\midrule
V3 suelo$\rightarrow$mesa & 1060,80 & 0 & 0 & 53,00 & 1,39 & 19,34 & 0,000 & -0,822 & 0,269 & 0,170 \\
V3 mesa$\rightarrow$suelo & 1060,34 & 0 & 0 & 68,00 & 1,93 & 14,67 & 0,206 & 0,143 & 0,374 & 0,233 \\
\bottomrule
\end{tabular}%
}
\caption{M\'etricas principales de captura y movimiento sobre los JSON V3.}
\label{tab:v3-current}
\end{table}

\begin{table}[h!]
\centering
\small
\resizebox{\textwidth}{!}{%
\begin{tabular}{lrrrrrr}
\toprule
Archivo & Idle ini. [s] & Idle fin. [s] & Tail recortada [s] & Fin activo Z [m] & Fin asentado Z [m] & Vel. residual [m/s] \\
\midrule
V3 suelo$\rightarrow$mesa & 2,37 & 1,77 & 0,63 & -1,684 & -0,822 & 0,000 \\
V3 mesa$\rightarrow$suelo & 2,54 & 2,34 & 0,33 & -0,394 & 0,143 & 0,000 \\
\bottomrule
\end{tabular}%
}
\caption{Separaci\'on entre endpoint activo y posici\'on asentada final tras estabilizaci\'on.}
\end{table}

\textbf{Lectura de los resultados:}
\begin{itemize}[leftmargin=14pt]
  \item El criterio de captura queda despejado: \textbf{BLE no domina} porque ambos archivos tienen \texttt{missingPackets = 0}.
  \item El cierre del tramo activo mejora visualmente: la trayectoria deja de prolongarse artificialmente durante la cola quieta y el residual reportado queda estabilizado a cero.
  \item La diferencia entre \textbf{fin activo} y \textbf{fin asentado} sigue siendo grande. Eso evidencia que la reconstrucci\'on de trayectoria a\'un arrastra error f\'isico dentro del propio movimiento, aunque la visualizaci\'on ya no perpet\'ua la deriva en reposo.
  \item El caso \emph{mesa $\rightarrow$ suelo} sigue preservando la forma ``sube un poco antes de bajar'', que era una se\~nal real del gesto y deb\'ia mantenerse.
\end{itemize}

\section{Por qu\'e el dibujo sale desvirtuado}
La deformaci\'on actual ya no se explica principalmente por la radio. Las causas m\'as probables son:
\begin{enumerate}[leftmargin=16pt]
  \item \textbf{Cierre tard\'io del movimiento}: aunque ya se recorta la cola quieta, el endpoint activo a\'un llega con mucha deriva acumulada.
  \item \textbf{Integraci\'on con velocidad residual}: el integrador sigue acumulando peque\~nos errores de aceleraci\'on/orientaci\'on y los convierte en un gran desplazamiento ficticio.
  \item \textbf{Sensibilidad del yaw/lateral}: sin referencia absoluta de heading, la lateral es especialmente sensible cuando el sensor gira o la manga de la barra rota.
  \item \textbf{Fijaci\'on de `barAxis'}: ahora es m\'as robusta, pero el sistema sigue siendo fr\'agil si el giro \'util no aporta suficiente observabilidad.
  \item \textbf{Sesgo temporal leve}: el reloj medio por muestra no parece roto, pero la combinaci\'on de \texttt{ts\_raw}, interpolaci\'on y distribuci\'on de \(\Delta t\) puede seguir introduciendo peque\~no sesgo cinem\'atico.
\end{enumerate}

\section{Hasta d\'onde puede mejorar IMU-only}
\textbf{Lo que s\'i esperamos mejorar todav\'ia:}
\begin{itemize}[leftmargin=14pt]
  \item detecci\'on m\'as limpia del inicio y del fin del gesto;
  \item conservaci\'on de la forma vertical del movimiento;
  \item altura m\'axima y altura asentada m\'as coherentes en gestos estructurados;
  \item lateral final menos contaminada por la cola de reposo;
  \item mayor consistencia entre gr\'aficas, tablas y visualizaci\'on 3D.
\end{itemize}

\textbf{Lo que no debemos prometer con esta configuraci\'on:}
\begin{itemize}[leftmargin=14pt]
  \item trayectoria 3D absoluta fiable;
  \item yaw absoluto estable;
  \item lateral precisa durante giro de manga o cambios de orientaci\'on sostenidos;
  \item un ``registro de vuelo'' tipo dron solo con esta IMU 6DoF.
\end{itemize}

En otras palabras: el techo realista de esta fase es una \textbf{buena descripci\'on del gesto}, no una navegaci\'on absoluta.

\section{Lecciones del paper}
El paper de referencia, \emph{A Study of a GNSS/IMU System for Object Localization and Spatial Position Estimation} (\textit{Sensors} 2025, 25, 6968, doi: \href{https://doi.org/10.3390/s25226968}{10.3390/s25226968}), es valioso como marco metodol\'ogico, pero no como benchmark directo para este proyecto.

Las ideas transferibles m\'as relevantes son:
\begin{itemize}[leftmargin=14pt]
  \item separar \textbf{actitud} y \textbf{cinem\'atica} en filtros distintos;
  \item modelar \textbf{bias} y ruido de sensores de forma expl\'icita;
  \item ajustar el filtro con \textbf{Allan variance} y par\'ametros estad\'isticos reales de las IMU;
  \item usar referencias absolutas (GNSS, magnet\'ometro, etc.) para corregir heading y posici\'on.
\end{itemize}

La diferencia fundamental es de observabilidad:
\begin{itemize}[leftmargin=14pt]
  \item el paper combina \textbf{GNSS + IMU 10DoF + magnet\'ometro + dos filtros de Kalman};
  \item nuestro caso actual es \textbf{IMU-only 6DoF} en un entorno con giro de manga, sin referencia absoluta de yaw ni de posici\'on.
\end{itemize}

Por eso el paper confirma la direcci\'on t\'ecnica correcta, pero tambi\'en deja claro por qu\'e no podemos exigirle a la fase actual una trayectoria absoluta de calidad comparable.

\section{Siguiente iteraci\'on}
La Iteraci\'on 2 deber\'ia centrarse en aumentar la observabilidad del sistema. La propuesta t\'ecnica m\'as coherente es:
\begin{enumerate}[leftmargin=16pt]
  \item \textbf{magnet\'ometro} bien calibrado para estabilizar yaw;
  \item posible \textbf{inclin\'ometro/p\'endulo MEMS} para robustecer actitud en escenarios con giro r\'apido;
  \item o un sensor auxiliar de \textbf{altura/posici\'on} si se quiere acercar el resultado a un ``flight log'';
  \item desacoplar de forma formal:
    \begin{itemize}[leftmargin=14pt]
      \item filtro de actitud;
      \item filtro cinem\'atico.
    \end{itemize}
\end{enumerate}

\section*{Conclusi\'on}
Los V3 son una buena noticia para el proyecto: demuestran que la captura BLE ya no est\'a rompiendo la sesi\'on. Eso nos permite atacar el problema real con m\'as precisi\'on: la reconstrucci\'on del movimiento con una IMU 6DoF, sin referencias absolutas, tiene un techo f\'isico claro.

La mejora implementada en esta iteraci\'on ---recorte de cola quieta, separaci\'on entre endpoint activo y endpoint asentado, y diagn\'osticos de cierre--- hace el sistema m\'as honesto y m\'as \'util. Describe mejor el gesto y deja visible cu\'ando la trayectoria ya entra en zona de baja confianza. Ese es exactamente el punto correcto para pasar a la siguiente iteraci\'on de arquitectura.

\end{document}
